\documentclass[11pt,fontset=fandol]{ctexart}
\usepackage[a4paper,margin=23mm]{geometry}
\usepackage{amsmath,amssymb,mathtools,bm,mathrsfs,esint}
\usepackage{graphicx,longtable,booktabs,array,calc}
\usepackage[normalem]{ulem}
\usepackage{xcolor,hyperref}
\hypersetup{hidelinks}
\setcounter{secnumdepth}{-1}
\newcounter{none}
\providecommand{\tightlist}{\setlength{\itemsep}{0pt}\setlength{\parskip}{0pt}}
\providecommand{\passthrough}[1]{#1}
\usepackage{amssymb,mathrsfs,esint,bm,newunicodechar}
\allowdisplaybreaks
\setlength{\emergencystretch}{3em}
\providecommand{\llangle}{\langle\!\langle}
\providecommand{\rrangle}{\rangle\!\rangle}
\newunicodechar{∞}{\ensuremath{\infty}}
\newunicodechar{∈}{\ensuremath{\in}}
\newunicodechar{∑}{\ensuremath{\sum}}
\newunicodechar{∫}{\ensuremath{\int}}
\newunicodechar{≤}{\ensuremath{\leq}}
\newunicodechar{≥}{\ensuremath{\geq}}
\newunicodechar{≠}{\ensuremath{\neq}}
\newunicodechar{→}{\ensuremath{\to}}
\newunicodechar{⇒}{\ensuremath{\Rightarrow}}
\newunicodechar{α}{\ensuremath{\alpha}}
\newunicodechar{β}{\ensuremath{\beta}}
\newunicodechar{γ}{\ensuremath{\gamma}}
\newunicodechar{λ}{\ensuremath{\lambda}}
\newunicodechar{φ}{\ensuremath{\varphi}}
\newunicodechar{ε}{\ensuremath{\varepsilon}}
\newunicodechar{σ}{\ensuremath{\sigma}}
\newunicodechar{θ}{\ensuremath{\theta}}
\newunicodechar{δ}{\ensuremath{\delta}}
\newunicodechar{η}{\ensuremath{\eta}}
\newunicodechar{π}{\ensuremath{\pi}}
\newunicodechar{ξ}{\ensuremath{\xi}}
\newunicodechar{Σ}{\ensuremath{\Sigma}}
\newunicodechar{Δ}{\ensuremath{\Delta}}
\newunicodechar{∂}{\ensuremath{\partial}}
\newunicodechar{∀}{\ensuremath{\forall}}
\newunicodechar{∃}{\ensuremath{\exists}}
\newunicodechar{√}{\ensuremath{\surd}}
\newunicodechar{−}{\ensuremath{-}}
\newunicodechar{μ}{\ensuremath{\mu}}
\newunicodechar{ρ}{\ensuremath{\rho}}
\newunicodechar{τ}{\ensuremath{\tau}}
\newunicodechar{ω}{\ensuremath{\omega}}
\newunicodechar{Ω}{\ensuremath{\Omega}}
\newunicodechar{χ}{\ensuremath{\chi}}
\newunicodechar{ψ}{\ensuremath{\psi}}
\newunicodechar{Γ}{\ensuremath{\Gamma}}
\newunicodechar{Φ}{\ensuremath{\Phi}}
\newunicodechar{Λ}{\ensuremath{\Lambda}}
\newunicodechar{ℝ}{\ensuremath{\mathbb{R}}}
\newunicodechar{ℤ}{\ensuremath{\mathbb{Z}}}
\newunicodechar{ℕ}{\ensuremath{\mathbb{N}}}
\newunicodechar{△}{\ensuremath{\triangle}}
\newunicodechar{⊗}{\ensuremath{\otimes}}
\begin{document}
\section{2024年杭州电子科技大学601数学分析补充解析}\label{ux5e74ux676dux5ddeux7535ux5b50ux79d1ux6280ux5927ux5b66601ux6570ux5b66ux5206ux6790ux8865ux5145ux89e3ux6790}

本文件由助手按回忆版题面独立编写，保留原章节、题号，供缺少解析的部分使用。原资料另附的解析涵盖一1---5、一8、二1---2、三1；其中二2的和函数推导不完整。本文件补全全部19个题号。题面明确缺失的一6、一7、三2仅说明缺失内容与可用方法，不补造题设。

\subsection{一、填空题}\label{ux4e00ux586bux7a7aux9898}

\subsubsection{1.
指数与幂的极限}\label{ux6307ux6570ux4e0eux5e42ux7684ux6781ux9650}

\textbf{答案：}\(2\)。

因为 \(n^{2024}/2^n\to0\)， \[
\log (n^{2024}+2^n)^{1/n}
=\log2+\frac1n\log\left(1+\frac{n^{2024}}{2^n}\right)\to\log2.
\] 对两端取指数即得答案。

\subsubsection{2. 不定积分}\label{ux4e0dux5b9aux79efux5206}

\textbf{答案：} \[
\int\frac{x^3}{\sqrt{1-x^2}}\,dx
=\frac13(1-x^2)^{3/2}-\sqrt{1-x^2}+C,\qquad |x|<1.
\] 令 \(u=1-x^2\)，则 \(x^3dx=-(1-u)du/2\)，所以积分等于 \[
-\frac12\int(u^{-1/2}-u^{1/2})\,du=-\sqrt u+\frac13u^{3/2}+C.
\]

\subsubsection{3.
麦克劳林级数系数}\label{ux9ea6ux514bux52b3ux6797ux7ea7ux6570ux7cfbux6570}

\textbf{答案：}\(\displaystyle\frac{(\log2)^{2024}}{2024!}\)。

利用 \[
2^x=e^{x\log2}=\sum_{k=0}^\infty\frac{(\log2)^k}{k!}x^k
\] 直接读取 \(x^{2024}\) 的系数。

\subsubsection{4.
极限函数的间断点}\label{ux6781ux9650ux51fdux6570ux7684ux95f4ux65adux70b9}

\textbf{答案：}\(x=0\)，为可去间断点。

固定 \(x\ne0\) 后将分子、分母除以 \(n\)，得到 \[
F(x)=\lim_{n\to\infty}\frac{(n+1)\sin x}{nx+4}
=\begin{cases}\sin x/x,&x\ne0,\\0,&x=0.\end{cases}
\] 因为 \(\lim_{x\to0}F(x)=1\ne F(0)\)，只有原点不连续。对某些负数
\(x\)，有限个下标可能使分母为零；上述极限按最终有定义的数列尾部理解，不影响极限函数。

\subsubsection{5.
两曲面交线的切线}\label{ux4e24ux66f2ux9762ux4ea4ux7ebfux7684ux5207ux7ebf}

\textbf{答案：}在交线任意一点 \((x_0,y_0,2)\)、\(x_0^2+y_0^2=2\)
处，切线为 \[
\boxed{\begin{cases}x_0x+y_0y=2,\\z=2.\end{cases}}
\] 令 \(r^2=x^2+y^2=z\ge0\)，由球面方程得到 \(z^2+z-6=0\)，故
\(z=2\)，交线是 \(z=2\) 平面上的圆 \(x^2+y^2=2\)。该圆在 \((x_0,y_0)\)
的切线即上式；等价的参数式为 \[
(x,y,z)=(x_0,y_0,2)+s(-y_0,x_0,0),\quad s\in\mathbb R.
\] 题面未指定切点，所以只能给出这一族切线。

\subsubsection{6.
题目未回忆到}\label{ux9898ux76eeux672aux56deux5fc6ux5230}

原卷仅写``没回忆到''，缺少题干、条件和待求量，无法确定题型或答案；也没有足够信息给出针对本题的解法。

\subsubsection{7.
方向导数，函数表达式缺失}\label{ux65b9ux5411ux5bfcux6570ux51fdux6570ux8868ux8fbeux5f0fux7f3aux5931}

已知 \(P_0=(-1,0,2)\)，\(P_1=(2,1,1)\)，方向单位向量是 \[
\boldsymbol\ell=\frac{(3,1,-1)}{\sqrt{11}}.
\] 若完整的 \(f\) 在 \(P_0\) 可微，则 \[
D_{\boldsymbol\ell}f(P_0)
=\frac{3f_x(P_0)+f_y(P_0)-f_z(P_0)}{\sqrt{11}}.
\] 原卷函数只写成
\(x^2+xy+\cdots\)，省略项的偏导数可能影响结果，不能据此前两项给出数值答案。

\subsubsection{8.
第一型曲线积分}\label{ux7b2cux4e00ux578bux66f2ux7ebfux79efux5206}

\textbf{答案：}\(\pi\)。

在单位上半圆上 \(x^2+y^2=1\)，故原积分等于上半圆的弧长，即 \(\pi\)。

\subsection{二、计算题}\label{ux4e8cux8ba1ux7b97ux9898}

\subsubsection{1. 极限}\label{ux6781ux9650}

\textbf{答案：}\(0\)。

由泰勒展开， \[
x-\sin x=\frac{x^3}{6}+O(x^5),\qquad
x(e^x-1)=x^2+O(x^3),
\] 所以原商等于 \(x/6+O(x^2)\)，极限为零。

\subsubsection{2.
幂级数的收敛域与和函数}\label{ux5e42ux7ea7ux6570ux7684ux6536ux655bux57dfux4e0eux548cux51fdux6570}

\textbf{答案：}收敛域 \((-1,1)\)，和函数为 \[
\boxed{-\frac{2x^2}{1+x^2}-\log(1+x^2)}.
\] 系数 \((2n+1)/n\to2\)，因此以 \(x\) 为变量的收敛半径为 \(1\)。在
\(x=\pm1\) 时通项不趋于零，均发散。在 \(|x|<1\) 内， \[
\begin{aligned}
\sum_{n=1}^\infty(-1)^n\frac{2n+1}{n}x^{2n}
&=2\sum_{n=1}^\infty(-x^2)^n+\sum_{n=1}^\infty\frac{(-x^2)^n}{n}\\
&=-\frac{2x^2}{1+x^2}-\log(1+x^2).
\end{aligned}
\]

\subsubsection{3.
上半圆上的第二型曲线积分}\label{ux4e0aux534aux5706ux4e0aux7684ux7b2cux4e8cux578bux66f2ux7ebfux79efux5206}

\textbf{答案：}\(\displaystyle\sin1+\frac\pi{16}\)。

注意 \[
(-e^y\cos x-xy)\,dx-e^y\sin x\,dy
=-d(e^y\sin x)-xy\,dx.
\] 从 \((1,0)\) 到 \((0,0)\) 的全微分积分是 \(\sin1\)。其余部分取参数 \[
x=\frac{1+\cos t}{2},\quad y=\frac{\sin t}{2},\quad0\le t\le\pi,
\] 则 \[
-\int_Lxy\,dx
=\frac18\int_0^\pi(1+\cos t)\sin^2t\,dt=\frac\pi{16}.
\] 两部分相加即得答案。

\subsection{三、证明题}\label{ux4e09ux8bc1ux660eux9898}

\subsubsection{\texorpdfstring{1. \(f(x,y)=|x-y|\varphi(x,y)\)
的偏导与可微性}{1. f(x,y)=\textbar x-y\textbar\textbackslash varphi(x,y) 的偏导与可微性}}\label{fxyx-yvarphixy-ux7684ux504fux5bfcux4e0eux53efux5faeux6027}

\textbf{结论：}两偏导数存在以及 \(f\) 在原点可微，均当且仅当
\(\varphi(0,0)=0\)；此时 \(f_x(0,0)=f_y(0,0)=0\)，\(df(0,0)=0\)。

令 \(c=\varphi(0,0)\)。因 \(f(0,0)=0\)，沿 \(x\) 轴的差商为 \[
\frac{f(h,0)}h=\operatorname{sgn}(h)\varphi(h,0),
\] 其右、左极限分别是 \(c,-c\)，故 \(f_x(0,0)\) 存在当且仅当 \(c=0\)。沿
\(y\) 轴同理，\(f_y(0,0)\) 也当且仅当 \(c=0\) 时存在。

当 \(c=0\) 时，由连续性及 \(|x-y|\le\sqrt2\sqrt{x^2+y^2}\)， \[
\frac{|f(x,y)|}{\sqrt{x^2+y^2}}
\le\sqrt2|\varphi(x,y)|\to0,
\] 所以 \(f\) 可微且微分为零。当 \(c\ne0\) 时偏导不存在，因而不可微。

\subsubsection{2.
第二型曲面积分，系数缺失}\label{ux7b2cux4e8cux578bux66f2ux9762ux79efux5206ux7cfbux6570ux7f3aux5931}

原题三个积分系数均未回忆到，无法计算数值。仅按已知的图形 \[
\Sigma:\ z=5-x^2-y^2,\quad x\ge1,
\] 以及``取上侧''，若积分式为
\(P\,dy\,dz+Q\,dz\,dx+R\,dx\,dy\)，则可化为 \[
\iint_D[2xP(x,y,z)+2yQ(x,y,z)+R(x,y,z)]_{z=5-x^2-y^2}\,dx\,dy,
\quad D=\{(x,y):x\ge1\}.
\] 这里 \(D\)
无界，除系数缺失外，还须确认原题是否另有边界；若没有，需要检验反常积分收敛性。不能自行添加
\(z\ge0\) 等条件。

\subsubsection{3.
含参量反常积分}\label{ux542bux53c2ux91cfux53cdux5e38ux79efux5206}

\textbf{答案：}\(I(x)=\arctan x\)，\(x\in\mathbb R\)。

由 \(|\sin(tx)|\le|tx|\)，被积函数在 \(t=0\)
附近可积，在无穷远处也绝对可积。对参数求导时 \[
\left|\frac\partial{\partial x}\left(e^{-t}\frac{\sin(tx)}t\right)\right|
=|e^{-t}\cos(tx)|\le e^{-t},
\] 故可由控制收敛定理交换求导与积分，得到 \[
I'(x)=\int_0^\infty e^{-t}\cos(tx)\,dt=\frac1{1+x^2}.
\] 最后一个积分可对 \(e^{-t}\cos(tx)\)、\(e^{-t}\sin(tx)\)
各分部积分一次求得。再由 \(I(0)=0\)，得 \(I(x)=\arctan x\)。

\subsection{四、综合题}\label{ux56dbux7efcux5408ux9898}

\subsubsection{\texorpdfstring{1. 用 \(\varepsilon\)---\(\delta\)
语言证明立方根保极限}{1. 用 \textbackslash varepsilon---\textbackslash delta 语言证明立方根保极限}}\label{ux7528-varepsilondelta-ux8bedux8a00ux8bc1ux660eux7acbux65b9ux6839ux4fddux6781ux9650}

对任意实数 \(u,v\)， \[
u^2+uv+v^2\ge\frac14(u-v)^2,
\] 所以 \(|u^3-v^3|\ge|u-v|^3/4\)。给定
\(\varepsilon>0\)，由右极限假设，存在 \(\delta>0\)，使
\(0<x-x_0<\delta\) 时 \[
|f(x)-A|<\varepsilon^3/4.
\] 令 \(u=\sqrt[3]{f(x)}\)，\(v=\sqrt[3]A\)，便有 \[
|\sqrt[3]{f(x)}-\sqrt[3]A|
\le(4|f(x)-A|)^{1/3}<\varepsilon.
\] 这正是所需右极限的定义，且同时覆盖 \(A=0\) 和负数 \(A\)。

\subsubsection{2.
一致连续与不一致连续}\label{ux4e00ux81f4ux8fdeux7eedux4e0eux4e0dux4e00ux81f4ux8fdeux7eed}

题面识读说明：远程OCR将第二个函数写作 \(g(x)=x+\sin x^2\)，本地OCR写作
\(x+\sin^2x\)，两份识读结果存在分歧。本节按远程题面的
\(g(x)=x+\sin(x^2)\) 作答；若第二式实际也是
\(x+\sin^2x\)，则``不一致连续''这一命题不成立。

对 \(f(x)=x+\sin^2x\)，\(f'(x)=1+\sin2x\)，故
\(|f'(x)|\le2\)。由中值定理， \[
|f(x)-f(y)|\le2|x-y|,
\] 所以 \(f\) 在整个实数轴上一致连续。

对 \(g(x)=x+\sin(x^2)\)，取 \[
x_n=\sqrt{2\pi n+\pi/2},\qquad y_n=\sqrt{2\pi n+3\pi/2}.
\] 有 \(y_n-x_n=\pi/(x_n+y_n)\to0\)，但 \[
g(y_n)-g(x_n)=(y_n-x_n)-2\to-2.
\] 这与一致连续的序列判据矛盾，故 \(g\) 不一致连续。

\subsubsection{\texorpdfstring{3. 证明某点
\(|f''|\ge4\)}{3. 证明某点 \textbar f\textquotesingle\textquotesingle\textbar\textbackslash ge4}}\label{ux8bc1ux660eux67d0ux70b9-fge4}

反设对所有 \(x\in(0,1)\) 都有 \(|f''(x)|<4\)。分别在 \([0,x]\) 和
\([x,1]\) 对 \(f'\) 用中值定理，结合 \(f'(0)=f'(1)=0\)，得到 \[
|f'(x)|<4x,\qquad |f'(x)|<4(1-x).
\] 由于 \(f'\) 连续且在内点严格小于上述界， \[
1=f(1)-f(0)=\int_0^1f'(x)\,dx
\le\int_0^1|f'(x)|\,dx
<4\int_0^1\min(x,1-x)\,dx=1,
\] 矛盾。因此存在 \(\xi\in(0,1)\) 使 \(|f''(\xi)|\ge4\)。此证法不要求
\(f''\) 连续。

\subsubsection{4.
积分与导数最大值的估计}\label{ux79efux5206ux4e0eux5bfcux6570ux6700ux5927ux503cux7684ux4f30ux8ba1}

令 \(M=\max_{[0,1]}|f'|\)。由端点值及中值定理， \[
|f(x)|\le Mx,\qquad |f(x)|\le M(1-x).
\] 因此 \[
\left|\int_0^1f(x)\,dx\right|
\le\int_0^1|f(x)|\,dx
\le M\left(\int_0^{1/2}x\,dx+\int_{1/2}^1(1-x)\,dx\right)=\frac M4.
\]

\subsubsection{5.
函数项级数的连续、可导与端点不可导}\label{ux51fdux6570ux9879ux7ea7ux6570ux7684ux8fdeux7eedux53efux5bfcux4e0eux7aefux70b9ux4e0dux53efux5bfc}

设 \[
f(x)=\sum_{n=1}^\infty\frac{e^{-nx}}{n^2+1},\quad x\ge0.
\]

\textbf{（1）在 \([0,\infty)\) 连续。}每项连续，且
\(|e^{-nx}/(n^2+1)|\le1/(n^2+1)\)。右侧正项级数收敛，故由 Weierstrass
判别法原级数在整个 \([0,\infty)\) 一致收敛，其和连续。

\textbf{（2）在 \((0,\infty)\) 可导。}任取 \(a>0\)，导数级数在
\([a,\infty)\) 上满足 \[
\left|\frac{-ne^{-nx}}{n^2+1}\right|\le e^{-na},
\] 所以一致收敛，可以逐项求导： \[
f'(x)=-\sum_{n=1}^\infty\frac{ne^{-nx}}{n^2+1},\qquad x>0.
\]

\textbf{（3）在 \(x=0\) 没有有限右导数。}对
\(h>0\)，每个差商项均非正，所以任意固定 \(N\) 都有 \[
\frac{f(h)-f(0)}h
=\sum_{n=1}^\infty\frac{e^{-nh}-1}{h(n^2+1)}
\le\sum_{n=1}^N\frac{e^{-nh}-1}{h(n^2+1)}.
\] 固定 \(h\) 时等式成立，因为两个函数值的级数均绝对收敛。令
\(h\downarrow0\)，右侧趋于 \(-\sum_{n=1}^N n/(n^2+1)\)。该部分和随
\(N\to\infty\) 趋于 \(-\infty\)，因此右差商趋于
\(-\infty\)。原函数只定义在
\([0,\infty)\)，这里的``不可导''具体指不存在有限右导数。

\subsection{核验说明}\label{ux6838ux9a8cux8bf4ux660e}

题面已交叉读取远程OCR与本地逐页OCR。使用 SymPy
对不定积分导数、计算题一的极限、计算题二的级数系数及计算题三的参数积分作符号检查；检查脚本为
\texttt{.agent/hdu-ocr/check\_math2024\_2025.py}。缺失题目与有识读歧义的条件均已在对应题号标记。
\end{document}
